\section{\texorpdfstring{\(L1\)}{L1}: a Zeckendorf protocol}\label{sec:L1}

Let \(F_1=1,F_2=2,F_{k+2}=F_{k+1}+F_k\).  Let
\(X_\infty^{\mathrm{fin}}\) be the finite-support binary sequences with no
adjacent \(1\)'s, and set
\(\Val(c)=\sum_{k\ge1}c_kF_k\).

\subsection{Normal forms}\label{ssec:fold}

\begin{definition}[Finite Zeckendorf interface]\label{def:fold}
For \(m\ge1\), put
\[
  \Omega_m=\{0,1\}^{m+1},
  \qquad
  X_m=\{c\in X_\infty^{\mathrm{fin}}:0\le \Val(c)<F_{m+2}\}.
\]
The finite valuation is
\[
  \Val_m\colon X_m\to \{0,1,\ldots,F_{m+2}-1\},
  \qquad \Val_m(c)=\Val(c).
\]
For a finite binary word \(\omega\in\Omega_m\), \(\Fold_m(\omega)\in X_m\) is
the greedy Zeckendorf normal form of the residue
\(\sum_{k=1}^{m+1}\omega_kF_k\pmod {F_{m+2}}\).  The stable normalizer
\(\mathrm{NF}_\infty\) is defined analogously for finite words of arbitrary
length, without reducing modulo a finite level.
\end{definition}

\begin{equation}\label{eq:fold-residue}
  \Val_m(\Fold_m(\omega))\equiv\sum_{k=1}^{m+1}\omega_kF_k\pmod {F_{m+2}} .
\end{equation}

\begin{proposition}[Finite Zeckendorf residue interface]
\label{prop:finite-zeckendorf-residue-interface}
The map \(\Val_m\) is a computable bijection
\[
  X_m\longrightarrow \ZZ/F_{m+2}\ZZ
\]
after identifying \(\ZZ/F_{m+2}\ZZ\) with its representatives
\(0,\ldots,F_{m+2}-1\).  Hence \(|X_m|=F_{m+2}\), and
\(\Fold_m\colon\Omega_m\to X_m\) is a computable surjection.
\end{proposition}

\begin{proof}
For each \(0\le n<F_{m+2}\), \Cref{lem:greedy-zeckendorf-normalizer} gives a
unique Zeckendorf representative \(c_n\in X_\infty^{\mathrm{fin}}\) with
\(\Val(c_n)=n\).  By definition this representative belongs to \(X_m\), so
\(\Val_m\) is surjective onto the displayed representative set.  If
\(c,d\in X_m\) and \(\Val_m(c)=\Val_m(d)\), then \(\Val(c)=\Val(d)\) as ordinary
integers in the interval \([0,F_{m+2})\), and Zeckendorf uniqueness gives
\(c=d\).  The greedy algorithm computes the inverse representative
\(n\mapsto c_n\), and reduction of \(\sum_{k=1}^{m+1}\omega_kF_k\) modulo
\(F_{m+2}\), followed by that inverse representative, computes \(\Fold_m\).  To
see surjectivity of \(\Fold_m\), let \(c\in X_m\) and let
\(\Val(c)=n<F_{m+2}\).  The Zeckendorf representative of such an \(n\) uses no
summand larger than \(F_{m+1}\); otherwise its value would be at least
\(F_{m+2}\).  Thus the first \(m+1\) digits of \(c\), viewed as an element of
\(\Omega_m\), fold back to \(c\).
\end{proof}

\begin{lemma}[Greedy Zeckendorf normalizer]
\label{lem:greedy-zeckendorf-normalizer}
Every \(n\in\NN\) has a unique representative \(c_n\in
X_\infty^{\mathrm{fin}}\), and the map \(n\mapsto c_n\) is computable.
\end{lemma}

\begin{proof}
This is the classical Zeckendorf theorem \cite{Zeckendorf1972}: repeatedly
subtract the largest Fibonacci number not exceeding the remaining integer.  The
greedy choice terminates, produces no adjacent summands, and uniqueness follows
from the standard exchange inequality
\(F_{k+1}>F_k+
F_{k-2}+F_{k-4}+\cdots\).
\end{proof}

\begin{proposition}[Computable Zeckendorf arithmetic anchor]
\label{prop:zeckendorf-computable-arithmetic-anchor}
The stable Zeckendorf normal-form set \(X_\infty^{\mathrm{fin}}\) gives a
computable presentation of \(\NN\): the valuation \(\Val\) and its greedy inverse
are computable, equality is decidable, and the operations transported from
ordinary addition and multiplication are computable on normal forms.  The only
finite-level operation used below is the residue fold of \Cref{def:fold}.  No
automatic-structure or Ostrowski-definability theorem is used as a hidden input
to the protocol theorem.
\end{proposition}

\begin{proof}
The map \(\Val\) is computed by the finite Fibonacci sum of the support, and
\Cref{lem:greedy-zeckendorf-normalizer}, the classical Zeckendorf theorem
\cite{Zeckendorf1972}, computes the inverse normal form.  Equality is decidable
by comparing finite binary strings after deleting trailing zeroes, or
equivalently by applying \(\Val\).  To compute transported addition or
multiplication of \(c,d\in X_\infty^{\mathrm{fin}}\), compute
\(a=\Val(c)\), \(b=\Val(d)\), form \(a+b\) or \(ab\) in ordinary integer
arithmetic, and apply the greedy inverse.  The finite fold is computable by
\Cref{prop:finite-zeckendorf-residue-interface}.  The broader literature gives
stronger finite-word algorithms and automatic/Ostrowski arithmetic interfaces
\cite{AhlbachEtAl2013,HieronymiTerry2018}; those results are cited as public
anchors for the numeration background, but the present argument needs only the
explicit computations just described.
\end{proof}

\begin{proposition}[Fold confluence]
\label{prop:fold-confluence}
Two finite words have the same stable fold if and only if they have the same
ordinary valuation.  At finite level \(m\), two words of \(\Omega_m\) have the
same finite fold if and only if their valuations have the same residue modulo
\(F_{m+2}\).  Hence equality of both stable and finite normal forms is decidable.
\end{proposition}

\begin{proof}
The stable assertion follows because both words fold to the unique Zeckendorf
representative of their ordinary valuation by
\Cref{lem:greedy-zeckendorf-normalizer}.  The finite assertion follows from the
bijection in \Cref{prop:finite-zeckendorf-residue-interface}.  Finite binary
words and residues are decidably comparable.
\end{proof}

\begin{definition}[Completions]\label{def:completions}
The arithmetic locus is
\begin{equation}\label{eq:X-infty}
  X_\infty^{\mathrm{fin}}
  =\{c\in\{0,1\}^{\NN}: c\text{ has finite support and no adjacent }1\text{'s}\}.
\end{equation}
The compact inverse-limit completion \(\widehat X\) is not used as an arithmetic
locus.
\end{definition}

\begin{remark}\label{rem:finite-vs-profinite}
The valuation is a bijection \(X_\infty^{\mathrm{fin}}\to\NN\), not a bijection
\(\widehat X\to\NN\).
\end{remark}

\subsection{Transported arithmetic}\label{ssec:arith}

\begin{definition}[Stable arithmetic]\label{def:stable-arith}
For \(c,d\in X_\infty^{\mathrm{fin}}\), define
\begin{align}
  c\oplus d&=\Val^{-1}(\Val(c)+\Val(d)),\label{eq:stable-add}\\
  c\otimes d&=\Val^{-1}(\Val(c)\Val(d)).\label{eq:stable-mul}
\end{align}
\end{definition}

\begin{proposition}[Finite residues are not stable arithmetic]
\label{prop:finite-residue-not-stable}
No fixed finite level \(X_m\) carries the stable arithmetic of
\((\NN,+,\times)\) through valuation.
\end{proposition}

\begin{proof}
The finite set \(X_m\) has only \(F_{m+2}\) elements, while addition or
multiplication on \(\NN\) leaves every finite interval.  Thus a fixed finite
residue level cannot be closed under the transported operations.
\end{proof}

\begin{proposition}[Finite residue-level transport]
\label{prop:finite-residue-level-transport}
For each \(m\), the bijection \(\Val_m\) transports the ring structure of
\(\ZZ/F_{m+2}\ZZ\) to \(X_m\).  Explicitly,
\[
\begin{aligned}
  c\oplus_m d&=\Val_m^{-1}(\Val_m(c)+\Val_m(d)\bmod F_{m+2}),\\
  c\otimes_m d&=\Val_m^{-1}(\Val_m(c)\Val_m(d)\bmod F_{m+2}).
\end{aligned}
\]
Then \(\Val_m\) is a finite ring isomorphism from
\((X_m,\oplus_m,\otimes_m)\) to \(\ZZ/F_{m+2}\ZZ\).  This is residue arithmetic
only and is not the stable arithmetic of \Cref{def:stable-arith}.
\end{proposition}

\begin{proof}
The displayed operations are the pullback of addition and multiplication in the
finite quotient ring through the bijection
\(\Val_m\colon X_m\to\ZZ/F_{m+2}\ZZ\) of
\Cref{prop:finite-zeckendorf-residue-interface}.  Therefore the ring axioms and
the isomorphism statement are immediate from the corresponding quotient-ring
facts.  The final sentence is \Cref{prop:finite-residue-not-stable}: a fixed
finite quotient records only residues modulo \(F_{m+2}\), while stable
arithmetic records ordinary natural-number values without quotienting.
\end{proof}

\begin{theorem}[Transported semiring]
\label{thm:semiring}
The map \(\Val\colon X_\infty^{\mathrm{fin}}\to\NN\) is an isomorphism from
\((X_\infty^{\mathrm{fin}},\oplus,\otimes)\) to \((\NN,+,\times)\).  At each
finite level, the only transported arithmetic asserted in this paper is the
residue ring interface of \Cref{prop:finite-residue-level-transport}.
\end{theorem}

\begin{proof}
This is immediate from the definitions of \(\oplus\) and \(\otimes\) and the
bijection in \Cref{lem:greedy-zeckendorf-normalizer}.  The finite-level clause
is \Cref{prop:finite-residue-level-transport}.
\end{proof}

Stable arithmetic in \Cref{thm:semiring} cannot be replaced by arithmetic on one
fixed finite residue set \(X_m\), by \Cref{prop:finite-residue-not-stable}.

\begin{lemma}[Value transport suffices]
\label{lem:value-transported-suffices}
Any computable bijection \(N\to\NN\) transports ordinary computable operations on
\(\NN\) to computable operations on \(N\).
\end{lemma}

\begin{proof}
Decode, apply the ordinary operation on \(\NN\), and encode back by the inverse
bijection.
\end{proof}

\begin{proposition}[Value-transport exact interface]
\label{prop:value-transport-exact-interface}
The Zeckendorf protocol uses only the computable bijection
\(\Val\colon X_\infty^{\mathrm{fin}}\to\NN\) and not any digit-local arithmetic
algorithm.
\end{proposition}

\begin{proof}
The operations are exactly those in \Cref{def:stable-arith}; both are obtained by
decode--compute--encode.
\end{proof}

\subsection{One primitive at two levels}\label{ssec:sole}

Let \(A\) be a nonempty set and let \(g\colon A\to A\) have infinite order.  Put
\[
  I_n(f)=f^{\circ n}\qquad(f\in\End(A)).
\]

\begin{theorem}[Two-level readout]
\label{thm:two-level}
With composition as the single semantic operation,
\begin{align}
  \Phi_0(c_n)&=g^{\circ n},\label{eq:level0}\\
  \Phi_1(c_n)&=I_n\label{eq:level1}
\end{align}
read addition and multiplication:
\begin{align}
  \Phi_0(c_{a+b})&=\Phi_0(c_a)\circ\Phi_0(c_b),\label{eq:add-from-comp}\\
  \Phi_1(c_{ab})&=\Phi_1(c_a)\circ\Phi_1(c_b).\label{eq:mul-from-comp}
\end{align}
\end{theorem}

\begin{proof}
The first identity is \(g^{a+b}=g^a\circ g^b\).  The second is
\(I_a\circ I_b(f)=(f^b)^a=f^{ab}=I_{ab}(f)\).
\end{proof}

Thus addition is read at the orbit level and multiplication at the power-operator
level; both readouts use composition as the only semantic operation.
\begin{lemma}[Level-zero exponent skeleton]
\label{lem:level-zero-exponent-skeleton}
Let \(C=\{g^n:n\in\NN\}\) be a free monogenic orbit.  Every multiplicative
readout \(\eta\colon\NN_{\ge1}\to C\) has the form
\[
  \eta(n)=g^{e(n)},
  \qquad e(ab)=e(a)+e(b),
\]
and hence
\[
  e(n)=\sum_{p}v_p(n)w(p)
\]
for uniquely determined weights \(w(p)=e(p)\in\NN\).
\end{lemma}

\begin{proof}
Since \(g
\) has infinite order, each element of \(C\) has a unique exponent.  The
homomorphism law for \(\eta\) becomes additivity of \(e\) on multiplication.
Unique prime factorization gives the displayed formula.
\end{proof}

\begin{theorem}[One-orbit multiplication obstruction and extremal skeleton]
\label{thm:level-zero-obstruction}
Fix one free monogenic orbit \(C=\{g^n:n\in\NN\}\) and require multiplication to
be read at level zero by the ambient product on \(C\).  Then every
multiplicative readout \(\eta\colon\NN_{\ge1}\to C\) is obtained from a unique
prime-weight function \(w\colon\mathcal P\to\NN\) by
\[
  \eta(n)=g^{e(n)},
  \qquad e(n)=\sum_p v_p(n)w(p),
\]
and no such readout is injective on all of \(\NN_{\ge1}\).  More sharply, if
\(H\subseteq\NN_{\ge1}\) is any multiplicative submonoid, then
\(\eta|_H\) is injective if and only if either \(H=\{1\}\), or there is an
integer \(q>1\) such that
\[
  H\subseteq\{q^k:k\in\NN\}
  \qquad\hbox{and}\qquad e(q)>0.
\]
Thus the obstruction is not a failure of computability or of normal forms: it is
the rank-one exponent skeleton forced by asking ambient product on one free
monogenic orbit to carry integer multiplication.  The positive replacement used
in the main hierarchy is the level shift to the Church power-operator readout of
\Cref{thm:two-level}.
\end{theorem}

\begin{proof}
The factorization and uniqueness of the weights are exactly
\Cref{lem:level-zero-exponent-skeleton}: the infinite order of \(g\) makes
exponents unique, the law \(\eta(ab)=\eta(a)\eta(b)\) makes \(e\) additive on
multiplication, and unique factorization gives
\(e(n)=\sum_p v_p(n)w(p)\) with \(w(p)=e(p)\).

This already excludes injectivity on all of \(\NN_{\ge1}\).  Choose distinct
primes \(p,r\).  If either \(w(p)=0\) or \(w(r)=0\), then the corresponding prime
has the same image as \(1\).  If both weights are positive, then
\[
  e(p^{w(r)})=w(p)w(r)=e(r^{w(p)}),
\]
so the two distinct prime powers have the same image.

It remains to prove the extremal statement for an arbitrary multiplicative
submonoid \(H\).  Suppose first that \(\eta|_H\) is injective.  If
\(h\in H\setminus\{1\}\) and \(e(h)=0\), then \(\eta(h)=g^0=\eta(1)\), a
contradiction.  Hence every nonunit in \(H\) has positive exponent weight.  For
\(a,b\in H\setminus\{1\}\), additivity gives
\[
  \eta(a^{e(b)})=g^{e(a)e(b)}=\eta(b^{e(a)}).
\]
Injectivity on \(H\) implies \(a^{e(b)}=b^{e(a)}\).  Comparing prime
factorizations yields
\[
  e(b)v_p(a)=e(a)v_p(b)
  \qquad\hbox{for every prime }p,
\]
so the exponent vectors of any two nonunit elements of \(H\) lie on the same
rational ray.  If \(H\ne\{1\}\), choose \(h\in H\setminus\{1\}\), let
\(\alpha=(\alpha_p)_p\) be the primitive nonnegative integer vector spanning the
rational ray through \((v_p(h))_p\), and set \(q=\prod_p p^{\alpha_p}\).  The
support is finite and \(q>1\).  Every \(x\in H\) then has exponent vector
\(k\alpha\) for some \(k\in\NN\), so \(x=q^k\).  Also \(e(q)>0\), because the
chosen element \(h=q^k\) is a nonunit of positive exponent weight.

Conversely, if \(H\subseteq\{q^k:k\in\NN\}\) and \(e(q)>0\), then
\(\eta(q^a)=\eta(q^b)\) implies \(ae(q)=be(q)\), hence \(a=b\).  Thus
\(\eta|_H\) is injective.  The final sentence is the positive identity
\(\Phi_1(c_{ab})=\Phi_1(c_a)\circ\Phi_1(c_b)\) from \Cref{thm:two-level}, which
reads multiplication after shifting from the orbit to the power-operator level.
\end{proof}

The obstruction applies only to same-level ambient-product readouts on one free
monogenic orbit.  It does not apply to the two-level readout of
\Cref{thm:two-level}, which reads multiplication by composition of power operators
in \(\End(\End(A))\), not by the ambient product on \(\{g^n\}\).  Within the
one-orbit model, \Cref{thm:level-zero-obstruction} is sharp: only the trivial
monoid or a single cyclic multiplicative ray can remain faithful.  The finer
prime-coordinate classification is split off from this short article.  No
finite-rank product-readout classification is used in this paper.

\begin{corollary}[Minimal two-level split]
\label{cor:minimal}
In the setting of \Cref{thm:two-level}, the two displayed readout levels are
jointly sufficient for the usual arithmetic semiring on \(\NN\): the map
\(n\mapsto c_n\) is injective, \(\Phi_0\) reads addition by composition in
\(\End(A)\), and \(\Phi_1\) reads multiplication by composition in
\(\End(\End(A))\).  Moreover, under the one-free-monogenic-orbit restriction in
which the same orbit \(C=\{g^{\circ n}:n\in\NN\}\) is required to carry an
injective multiplication readout by its ambient product, no level-zero-only
replacement can be injective on all of \(\NN_{\ge1}\).  Hence within that
restricted model the passage from the orbit level to the power-operator level is
the minimal extra readout needed to keep one semantic primitive, composition,
while reading both \(+\) and \(\times\).
\end{corollary}

\begin{proof}
The sufficiency assertions are exactly the identities of \Cref{thm:two-level}:
\(g^{\circ(a+b)}=g^{\circ a}\circ g^{\circ b}\) at level zero and
\(I_a\circ I_b=I_{ab}\) at the power-operator level.  The symbols \(c_n\) are
indexed by \(n\in\NN\), so \(n\mapsto c_n\) is injective as a syntactic normal
form map; when read at level zero, the infinite order of \(g\) also makes
\(n\mapsto g^{\circ n}\) injective.

It remains only to justify the stated minimality, which is restricted to the
one-orbit model named in the statement.  If multiplication were read there by a
level-zero map \(\eta\colon\NN_{\ge1}\to C\) satisfying
\(\eta(ab)=\eta(a)\eta(b)\) for the ambient product on \(C\), then unique
exponents in the free monogenic orbit would give
\(\eta(n)=g^{\circ e(n)}\) and \(e(ab)=e(a)+e(b)\).  This is the exponent skeleton
proved in \Cref{lem:level-zero-exponent-skeleton}.  The prime-power collision in
\Cref{thm:level-zero-obstruction} then shows that such an \(\eta\) cannot be
injective on all of \(\NN_{\ge1}\).  Thus a level-zero-only replacement fails
precisely in the restricted model, while \Cref{thm:two-level} supplies the next
readout level and realizes multiplication there by the same semantic operation,
composition.
\end{proof}


\begin{definition}[Protocol universality]\label{def:prot-univ}
An \(L1\) protocol consists of computable normal forms with decidable equality,
compiled operations on those normal forms, and semantic readout identities for
the operations being represented.
\end{definition}

In the Zeckendorf protocol, the compiled arithmetic operations on normal forms
and the semantic readouts \(\Phi_0,\Phi_1\) are separate pieces of data.  The
compiled operations are \(\oplus,\otimes\) from \Cref{def:stable-arith}; the
readouts are the maps in \Cref{thm:two-level}.  All computability in the \(L1\)
construction occurs on the finite-support normal forms
\(X_\infty^{\mathrm{fin}}\), by
\Cref{lem:greedy-zeckendorf-normalizer,prop:value-transport-exact-interface}.

\begin{theorem}[Restricted level-zero protocol dichotomy]
\label{thm:restricted-level-zero-protocol-dichotomy}
\label{cor:level-zero-protocol-dichotomy}
Let a protocol for multiplication on \(\NN_{\ge1}\) use a single free
monogenic orbit \(C=\{g^n:n\in\NN\}\) as its semantic multiplication target, and
suppose the multiplication readout is a map
\(\eta\colon\NN_{\ge1}\to C\) satisfying
\[
  \eta(ab)=\eta(a)\eta(b)
  \qquad(a,b\ge1),
\]
where the product on the right is the ambient monoid product on \(C\).  Thus
the assertion is made under this level-zero monogenic-orbit restriction.  Then
\(\eta\) is not injective.  Consequently any protocol that realizes
\(\{+,\times\}\) while satisfying this same level-zero one-free-monogenic-orbit
restriction must, for its multiplication readout, leave the one free monogenic
orbit, abandon ambient-product readout, or move to a higher readout level such
as the power-operator level of \Cref{thm:two-level}.  The theorem makes no
global rigidity assertion about protocols outside this restricted model.
\end{theorem}

\begin{proof}
The hypotheses are exactly the hypotheses of \Cref{thm:level-zero-obstruction},
so the readout \(\eta\) is not injective.  Hence an injective protocol cannot
simultaneously keep multiplication in one free monogenic orbit and read it by
the ambient product.  The displayed alternatives are the logical negations of
those two restrictions, together with the positive realized alternative:
\Cref{thm:two-level} reads multiplication by composition of the operators
\(I_n\) in \(\End(\End(A))\), not by the ambient product on \(C\).
\end{proof}

\begin{theorem}[Main-chain protocol rigidity boundary]
\label{prop:no-global-protocol-rigidity-input}
\label{thm:main-chain-protocol-rigidity-boundary}
The hierarchy theorem uses only the restricted dichotomy of
\Cref{thm:restricted-level-zero-protocol-dichotomy}.  It does not assert a
classification or rigidity theorem for all single-primitive protocols realizing
\(\{+,\times\}\).  Equivalently, every negative protocol inference in
\Cref{thm:main-hierarchy} factors through the one-free-monogenic-orbit hypothesis
and the ambient-product multiplication readout; protocols outside that
restricted model are judged only by the positive requirements of
\Cref{def:prot-univ} and by the explicit Zeckendorf realization of
\Cref{thm:zeckendorf-protocol-univ}.
\end{theorem}

\begin{proof}
The positive protocol is the value-transported Zeckendorf protocol of
\Cref{thm:zeckendorf-protocol-univ}; its multiplication readout is the
higher-level identity \(\Phi_1(c_{ab})=\Phi_1(c_a)\circ\Phi_1(c_b)\) from
\Cref{thm:two-level}.  The negative protocol input cited by
\Cref{thm:main-hierarchy} is \Cref{thm:level-zero-obstruction}, equivalently the
restricted dichotomy just proved.  Neither statement ranges over arbitrary
protocol data in \Cref{def:prot-univ}: the negative statement assumes a free
monogenic orbit \(\{g^n:n\in\NN\}\) and asks that multiplication be read by the
ambient product on that same orbit.  If either hypothesis is removed, the proof
has no exponent map \(e(ab)=e(a)+e(b)\) and hence no prime-weight skeleton from
\Cref{lem:level-zero-exponent-skeleton}.  Thus the only proved obstruction is
the restricted one, while the constructive Zeckendorf protocol supplies the
separate positive realization.  No global rigidity classification is therefore
used, proved, or needed in the main theorem chain.
\end{proof}

\begin{theorem}[Value-transported Zeckendorf protocol universality]
\label{thm:zeckendorf-protocol-univ}
The raw input grammar is the decidable set \(\{0,1\}^{<\infty}\) of finite
binary words.  The stable normalizer
\[
  \mathrm{NF}_\infty\colon \{0,1\}^{<\infty}\to X_\infty^{\mathrm{fin}}
\]
sends a word \(\omega=(\omega_1,\ldots,\omega_\ell)\) to the greedy Zeckendorf
normal form of \(\sum_{k=1}^\ell \omega_kF_k\).  With this raw-word normalizer,
the Zeckendorf normal forms, transported operations \(\oplus,\otimes\), and
readouts \(\Phi_0,\Phi_1\) form an \(L1\) protocol universal for
\(\{+,\times\}\) on \(\NN\).  It is a value-transported protocol and asserts no
bounded-window or digit-local arithmetic theorem.
\end{theorem}

\begin{proof}
The set \(\{0,1\}^{<\infty}\) is decidable, and the value
\(\sum_{k=1}^\ell\omega_kF_k\) of a raw word is computable.  Applying the greedy
Zeckendorf algorithm of \Cref{lem:greedy-zeckendorf-normalizer} to that integer
therefore computes \(\mathrm{NF}_\infty(\omega)\), and
\Cref{prop:fold-confluence} gives decidable equality of the resulting normal
forms.  The compiled operations on the image
\(X_\infty^{\mathrm{fin}}\) are computable by
\Cref{prop:zeckendorf-computable-arithmetic-anchor,lem:value-transported-suffices}.
The readout identities are
\Cref{thm:two-level}.  The final sentence is
\Cref{prop:value-transport-exact-interface}.
\end{proof}

The positive \(L1\) result uses no operation on \(\NN\) beyond ordinary addition
and multiplication transported through the Zeckendorf valuation; this is exactly
\Cref{def:stable-arith} and \Cref{thm:zeckendorf-protocol-univ}.

